\documentclass[a4paper,USenglish,anonymous]{lipics-v2021}
\bibliographystyle{plainurl}
\hideLIPIcs
\usepackage{booktabs}
\usepackage{stmaryrd}
\SetSymbolFont{stmry}{bold}{U}{stmry}{m}{n}
\newcommand{\BV}{\mathbb{B}_{32}}
\newcommand{\Fin}{\mathsf{Fin}}
\newcommand{\Ten}{\mathsf{Tensor}}
\newcommand{\Src}{\mathsf{S}}
\newcommand{\Tgt}{\mathsf{T}}
\newcommand{\lo}{\mathcal{L}}
\newcommand{\word}[1]{\langle #1\rangle_{32}}
\newcommand{\se}[1]{\llbracket #1\rrbracket^{\Src}_{\rho}}
\newcommand{\te}[1]{\llbracket #1\rrbracket^{\Tgt}_{\rho}}
\newcommand{\op}[1]{\ensuremath{\mathsf{#1}}}
\newcommand{\code}[1]{\texttt{\detokenize{#1}}}
\title{A Mechanized Tensor Core for TTIR-to-TTAdapter Lowering}
\titlerunning{A Mechanized Tensor Core for TTIR-to-TTAdapter Lowering}
\author{Anonymous Author}{}{}{}{}
\authorrunning{Anonymous Author}
\Copyright{Anonymous Author}
\ccsdesc[500]{Software and its engineering~Formal software verification}
\ccsdesc[500]{Software and its engineering~Compilers}
\keywords{tensor compilers, Triton, MLIR, Lean, semantic preservation, bit vectors}
\begin{document}
\maketitle
\begin{abstract}
Tensor compiler transformations often look elementary while depending on shape
conditions, finite-width arithmetic, and initialization conventions. We present
a source-audited Lean formalization of two conversion patterns in an Ascend
Triton compiler: splat and range construction. The source language contains
rank-one i32 tensors, inputs, splats, bounded ranges, and pointwise addition.
The target language separates index generation, fill, insertion into an
arbitrarily initialized tensor, and addition. A shape-indexed lowering selects
insertion for singleton splats and fill otherwise; range lowering either
returns generated indices or adds a filled signed offset. We prove equality
of source and target tensor values for every expression and environment, and
derive independence from the initial contents assigned to empty tensors.
Ten checked examples exercise signed starts, arithmetic wraparound, and two
incorrect hypothetical translations. The development uses exact 32-bit words
and includes a principal-theorem axiom audit. Its compiler correspondence is
manual: TTAdapter is a mixed-IR pipeline stage, and neither the full stage nor
the C++ implementation is verified. The resulting case study makes the
local semantic obligations precise while exposing the additional evidence
needed for implementation-level correctness.
\end{abstract}

\section{Introduction}
Triton expresses tiled computations through an intermediate language and
compiler designed for neural-network workloads~\cite{triton}. In an
MLIR-based implementation, a tensor operation can be represented at several
levels of abstraction before reaching executable code. MLIR deliberately
supports such heterogeneous representations and transformations~\cite{mlir}.
For verification, this flexibility raises a basic question: which semantic
claim is justified by a proof of a particular lowering pattern?

Consider a scalar splat into a tensor. Filling every lane with the scalar
appears straightforward, but a compiler can instead insert the scalar into
an empty tensor when the result has exactly one element. The optimization
depends on a shape fact: the single write must cover every observable lane.
Similarly, a range beginning at zero can be implemented by generated loop
indices, whereas a nonzero range requires an offset. The offset is signed,
the result is a 32-bit word, and addition is not unrestricted integer
addition. Ignoring any of these distinctions can make a superficially
plausible account of the translation inaccurate.

We formalize these two patterns from a pinned Ascend Triton checkout. The
destination artifact is called \code{ttadapter} by its Python compiler
driver. It is a pipeline stage, not the name of a single MLIR dialect. We
therefore introduce an abstract target language for the selected computations
rather than claiming to define the semantics of every operation in a
TTAdapter module. Our source language is similarly a restricted tensor core,
not the complete TTIR language.

The central theorem is a universal equality, not an enumeration of tests.
For every source expression, every shape permitted by its constructors, and
every environment, evaluating the lowered expression yields the same tensor
as evaluating the source expression. Both languages use 32-bit bit vectors.
The target explicitly distinguishes a full fill from insertion at index zero
into a tensor with arbitrary initial contents. Consequently, a corollary
states that the output of the lowering cannot depend on those contents.

This result has a deliberately limited interpretation. The source and target
denotations are manually specified. The translation is a Lean function
written to reflect inspected C++ branches, not code extracted from that
compiler. We neither execute the compiler nor compare generated MLIR in
this study. The contribution is a small, checked semantic specification and
an explicit implementation correspondence, not verification of Triton,
Ascend, or an end-to-end compilation pipeline.

\paragraph{Contributions.}
\begin{itemize}
\item A shape-indexed, rank-one i32 source and target core that distinguishes
signed range attributes, modular arithmetic, and partial initialization.
\item Mechanized preservation proofs for both splat branches, both range
branches, and their expression-tree compositions with addition.
\item A corollary proving independence from arbitrary empty-tensor contents,
and ten checked examples including two hypothetical bad-lowering witnesses.
\item A revision-specific source audit that separates proven abstractions
from unverified compiler, representation, and execution obligations.
\end{itemize}
We make no claim of a new general compiler-verification technique or priority
over all tensor formalizations. The value of the case study is its precise
scope, correspondence to concrete implementation branches, and reproducible
proof artifact.

\section{Compiler Boundary and Motivating Patterns}\label{sec:source}
The audited checkout is the \code{3rdparty/triton-ascend} subdirectory of a
Triton-distributed-ascend repository at revision
\begin{center}
\code{fdecef82c42cade68b4039af35a0e06ac8851536}.
\end{center}
All source anchors in this paper refer to that revision. We use the operation
declarations, verifier, conversion patterns, and Python driver as evidence
about implementation structure. This is a manual audit; the revision is not
a certificate relating executable compiler behavior to Lean.

\subsection{The stage is larger than the modeled transformation}
The driver registers \code{stages["ttadapter"]} through
\code{ttir_to_linalg} with named operations enabled. That function invokes
several passes, including control-flow and structure transformations, before
returning the transformed module text. Temporary and dumped output uses the
name \code{kernel.ttadapter.mlir}. A pure-SIMT route can bypass this stage.
Thus, even complete proofs of the selected conversion patterns would not
establish preservation by every preceding or following pass.

Our mathematical boundary is a pure tensor computation with fixed inputs.
We retain only the value-relevant parts of splat and range construction and
provide addition for composition. Source addition is an abstract signless
i32 arithmetic operation, not a claimed \code{tt.add} instruction. Inputs
allow the theorem to quantify over arbitrary tensors rather than only
closed constant expressions.

\subsection{Splat: a full write or a singleton write}
The inspected splat converter creates an empty result tensor. If every
dimension equals one, it inserts the scalar at all-zero indices. Otherwise
it emits a fill. For the rank-one shapes modeled here, the test reduces to
whether the length is one. It is incorrect to describe this implementation
as unconditionally lowering splat to fill.

For a singleton tensor, insertion and fill have the same value at the sole
valid index. For length four, inserting only at index zero does not establish
the values of the other three lanes. The formal target makes this distinction
observable by assigning arbitrary values to the initially empty tensor.
This is a mathematical model of unspecified contents, not a semantics of
poison propagation or an assertion that an actual allocation contains
particular machine words.

\subsection{Range: generated indices and a signed offset}
The range converter constructs a \code{linalg.generic} operation with
identity indexing and a parallel iterator. Its body obtains the loop index,
casts it to i32, and yields it. With start zero, this tensor is the result.
With nonzero start, the converter fills another tensor with the signed start
constant and combines it with the index tensor using \code{arith.addi}.

The declaration of \code{tt.make_range} explicitly warns that generated
accessors can return unsigned values although the attributes represent
signed i32. The inspected converter and verifier read the signed attribute
through \code{getStartAttr().getInt()}. Our model consequently encodes the
start with \code{BitVec.ofInt}, not an unchecked conversion to a natural
number. An example starting at $-2$ tests this choice.

The verifier requires a nonempty interval, a rank-one result, a length equal
to the difference between endpoints, and i32 elements. We encode a restricted
well-formed range directly in the source constructor. We do not model parsing
attributes or prove equivalence with the full verifier.

\section{A Shape-Indexed Tensor Core}\label{sec:core}
\subsection{Words, shapes, and environments}
Let $\BV$ be the set of 32-bit words, represented by Lean's
\code{BitVec 32}. Write $\word{z}$ for conversion of an integer $z$ to its
32-bit representation, and $u\oplus v$ for addition modulo $2^{32}$.
For a natural number $n$, define
\[
 \Fin(n)=\{i\in\mathbb N\mid i<n\},\qquad
 \Ten(n)=\Fin(n)\to\BV.
\]
The Lean representation is \code{Fin n → BitVec 32}. This is an extensional
tensor semantics: two tensors are equal when their values at every valid
index are equal. It says nothing about physical layout or memory addresses.

An environment $\rho=(I,U)$ contains two total families:
\[
 I : \prod_{n:\mathbb N}(\mathbb N\to\Ten(n)),
 \qquad U : \prod_{n:\mathbb N}\Ten(n).
\]
$I(n,x)$ supplies input identifier $x$ at shape $n$, while $U(n)$ supplies
the initial contents used by the target's insertion operation. Input
identifiers are shape-indexed; the same natural number at a different shape
does not impose a consistency condition. The model does not implement SSA
names, dominance, or a variable-binding discipline.

The family $U$ chooses one tensor per shape. It does not provide independent
fresh choices for each allocation, nor temporal nondeterminism. Our theorem
is independence from this explicitly encoded family. Establishing refinement
from a more detailed nondeterministic or poison-aware semantics remains a
separate obligation.

\subsection{The range domain}
For a signed start $s\in\mathbb Z$ and length $n\in\mathbb N$, set
\begin{equation}\label{eq:valid}
 V(s,n)\quad\equiv\quad
 0<n\;\land\;-2^{31}\leq s\;\land\;s+n\leq 2^{31}-1.
\end{equation}
The exclusive end is represented by $s+n$, not a separate stored parameter.
Since the interval is nonempty, these inequalities also bound the start
above and the end below. Both endpoints are representable as signed i32.
An exclusive end of $2^{31}$ is intentionally excluded even if some
mathematical lane sequence ending just before it would be representable.

The \code{makeRange} constructor carries a proof of $V(s,n)$. Range
validation is therefore a construction obligation rather than a possible
failure of the evaluator. Other constructors may have length zero: their
denotations are empty functions. This does not assert that every such shape
is legal in the actual TTIR implementation.

\subsection{Indexed syntax}
For each $n$, the source syntax $\Src_n$ has the following constructors:
\[
 e ::= \op{input}_n(x)\mid\op{splat}_n(v)
       \mid\op{range}_n(s,h)\mid e_1+_n e_2,
 \qquad h:V(s,n).
\]
Both operands of $+_n$ belong to $\Src_n$. The target syntax $\Tgt_n$ is
\[
 t ::= \op{input}_n(x)\mid\op{fill}_n(v)
       \mid\op{insert0}_n(v)\mid\op{index}_n\mid t_1+_n t_2.
\]
Here $v\in\BV$. The size indices are part of the inductive Lean types,
so ill-shaped additions cannot be constructed. In particular, a function
of type $\Src_n\to\Tgt_n$ preserves the represented length by construction.
This intrinsic property must not be confused with a separate theorem
about the real compiler's type converter.

\subsection{Source and target denotations}
For $i\in\Fin(n)$, source evaluation is defined by
\begin{align*}
 \se{\op{input}_n(x)}(i)&=I(n,x)(i),\\
 \se{\op{splat}_n(v)}(i)&=v,\\
 \se{\op{range}_n(s,h)}(i)&=\word{s}\oplus\word{i},\\
 \se{e_1+_n e_2}(i)&=\se{e_1}(i)\oplus\se{e_2}(i).
\end{align*}
The range denotation is modular from the outset. Bounds on the corresponding
mathematical integers are proved separately. We have not mechanized a general
bridge stating that signed decoding of every computed range word equals
$s+i$ under $V(s,n)$. Such a bridge is not needed by the preservation proof,
but would strengthen correspondence to a mathematical range specification.

Target evaluation uses the following equations:
\begin{align*}
 \te{\op{input}_n(x)}(i)&=I(n,x)(i),\\
 \te{\op{fill}_n(v)}(i)&=v,\\
 \te{\op{insert0}_n(v)}(i)&=
       \begin{cases}v&i=0,\\U(n)(i)&i\ne0,\end{cases}\\
 \te{\op{index}_n}(i)&=\word{i},\\
 \te{t_1+_n t_2}(i)&=\te{t_1}(i)\oplus\te{t_2}(i).
\end{align*}
The \op{index} constructor compresses index generation, cast, and yield
into one pure operation. It presupposes that the surrounding index
representation can express the modeled loop indices; no MLIR data-layout
or index-width theorem is provided. The \op{insert0} constructor has a
total denotation even at length zero, where there are no valid observations.
It is not a semantics of invalid \code{tensor.insert} operations. The
lowering emits it only at length one.

\section{Lowering and Semantic Preservation}\label{sec:proofs}
Define $\lo:\Src_n\to\Tgt_n$ structurally. Inputs are preserved and
addition lowers recursively. The interesting cases are
\begin{align}
 \lo(\op{splat}_n(v))&=
 \begin{cases}
 \op{insert0}_n(v)&n=1,\\
 \op{fill}_n(v)&n\ne1,
 \end{cases}\label{eq:splat}\\
 \lo(\op{range}_n(s,h))&=
 \begin{cases}
 \op{index}_n&s=0,\\
 \op{index}_n+_n\op{fill}_n(\word{s})&s\ne0.
 \end{cases}\label{eq:range}
\end{align}
The range witness $h$ does not affect the generated target expression.
It establishes membership in the chosen source domain. In fact, the
bit-vector identity underlying the range proof does not require its bounds;
the bounds concern admissibility and the intended signed-range reading.

\begin{lemma}[Splat preservation]\label{thm:splat}
For every $n,v,\rho$, evaluating the right-hand side of
Equation~\eqref{eq:splat} yields the constant tensor with value $v$.
\end{lemma}
\begin{proof}
By extensionality, fix $i\in\Fin(n)$. If $n=1$, then $0\leq i<1$
implies $i=0$. Target insertion therefore returns $v$, not $U(n)(i)$.
If $n\ne1$, the target is a fill and returns $v$ at every valid index.
These are exactly the source splat values. The length-zero case in the
second branch is vacuous.
\end{proof}
This is \code{lower_splat_correct}. Its singleton argument is not an
assumption that empty tensors are zero: it proves that no initial value is
observed. This distinction is essential to the negative example in
Section~\ref{sec:examples}.

\begin{lemma}[Range preservation]\label{thm:range}
For every $n,s,h:V(s,n),\rho$, evaluation of the right-hand side of
Equation~\eqref{eq:range} equals $i\mapsto\word{s}\oplus\word{i}$.
\end{lemma}
\begin{proof}
Fix a valid index $i$. If $s=0$, the target returns $\word{i}$ and the
source returns $\word{0}\oplus\word{i}$, equal by the zero identity.
Otherwise, the target returns $\word{i}\oplus\word{s}$. Commutativity
of bit-vector addition gives the source value. Extensionality gives equality
of the tensors.
\end{proof}
The Lean theorem \code{lower_range_correct} uses
\code{BitVec.add_comm}. It establishes the equality for finite-width
operations; no silent replacement by unbounded addition occurs.

\begin{theorem}[Lowering correctness]\label{thm:main}
For every $n$, environment $\rho$, and expression $e\in\Src_n$,
\[\te{\lo(e)}=\se{e}.\]
\end{theorem}
\begin{proof}
Proceed by structural induction on $e$. The input case follows from the
shared environment lookup. Lemmas~\ref{thm:splat} and~\ref{thm:range}
discharge the splat and range cases. For $e_1+_n e_2$, fix a lane $i$.
The target adds the values of $\lo(e_1)$ and $\lo(e_2)$ at $i$. The
induction hypotheses replace these with the source values of $e_1$ and
$e_2$, after which the source addition equation applies. Extensionality
finishes the case.
\end{proof}
The machine-checked statement is \code{lowering_correct}. It covers all
finite expression trees in the grammar, arbitrary input tensors, and
arbitrary $U$.

The corollary \code{lowering_lane_correct} specializes the equality to
any $i\in\Fin(n)$. There is no execution
trace, concurrency relation, or termination claim about the compiler here.

\subsection{Independence from initial contents}
Source evaluation never consults $U$. Structural induction proves that
environments with equal input families give equal source results; the
artifact calls this \code{source_ignores_empty}. Combining this property
with Theorem~\ref{thm:main} gives a target property.

\begin{corollary}[Empty-content independence]\label{thm:empty}
If $\rho_1=(I,U_1)$ and $\rho_2=(I,U_2)$, then for every $e\in\Src_n$,
\[
 \llbracket\lo(e)\rrbracket^{\Tgt}_{\rho_1}
 =\llbracket\lo(e)\rrbracket^{\Tgt}_{\rho_2}.
\]
\end{corollary}
\begin{proof}
Rewrite each side to its source denotation using Theorem~\ref{thm:main},
then apply source independence.
\end{proof}
This is \code{lowering_ignores_empty}. It is stronger than checking
results with a particular zero-initialized environment, but narrower than
proving absence of undefined behavior in generated MLIR. The theorem is
about the stated value model and the image of $\lo$, not every target term.
Indeed, a general \op{insert0} term at length greater than one may depend
on $U$.

\subsection{Mathematical range bounds}
The artifact separately checks
\begin{align*}
 s&\leq s+i<s+n,\\
 -2^{31}&\leq s+i<2^{31}-1
\end{align*}
for a valid range and lane. These are \code{range_lane_bounds} and
\code{range_lane_signed_bounds}. The first inequality only needs
$i\in\Fin(n)$; its range-validity parameter is unused in the proof.
The second uses the endpoint bounds. Its strict upper bound follows from
the exclusive end, itself bounded above by $2^{31}-1$. These integer facts
are not a substitute for the unproved general signed-decoding bridge
identified in Section~\ref{sec:core}.

\section{Mechanization and Checked Examples}\label{sec:examples}
The implementation uses Lean 4.19.0 and its bundled \code{Std} library,
with no external library dependencies. Lean supports both the functional
definitions and their proofs in one language~\cite{lean4}. The development
separates the core definitions, lowering proofs, examples, and an axiom
inspection file. It is independent of the neighboring synchronization
formalization; no synchronization result is imported to justify tensor
lowering.

\begin{table}[t]
\caption{Artifact structure and checked evidence. Paths are relative to the
formalization root, one directory above this paper.}\label{tab:artifact}
\centering
\begin{tabularx}{\linewidth}{@{}lX@{}}
\toprule
File & Role \\
\midrule
\code{TTIRFormal/Core.lean} & Indexed syntax, denotations, two integer-bound theorems.\\
\code{TTIRFormal/Lowering.lean} & Lowering, two local preservation results, composition, lane equality, and two independence results.\\
\code{TTIRFormal/Examples.lean} & Ten concrete example and counterexample theorems.\\
\code{TTIRFormal/Audit.lean} & Axiom dependencies of eight principal results.\\
\bottomrule
\end{tabularx}
\end{table}

\subsection{Positive cases}
The example environment supplies zero-valued inputs and assigns 99 to
every initially empty lane. A length-four range starting at $-2$ checks
signed decoding of its first lane as $-2$ and its last lane as 1. A
zero-start range checks its last lane as 3. A singleton splat checks the
insert branch with value 7. Adding a length-four splat of 7 to the negative
range checks the final lane as 8 and exercises composition and the fill
branch. Finally, adding singleton splats of $2^{32}-1$ and 1 checks a
result of zero, demonstrating that the source arithmetic is modular.

These six positive theorems complement rather than establish the universal
result. For example, the wraparound case checks the chosen arithmetic
interpretation, while the induction theorem already quantifies over every
input word. No benchmark timing or performance comparison is asserted.

\subsection{Two hypothetical incorrect translations}
The first counterexample replaces the nonzero-start range lowering with
bare index generation. At lane zero of the range beginning at $-2$, the
bad target produces $\word{0}$ and the source produces $\word{-2}$.
Their inequality is checked by \code{dropped_offset_counterexample}.
This is a witness against an alternative mathematical translation, not
an observed defect or mutation test run against the C++ compiler.

The second uses singleton insertion for a length-four splat of 7. At lane
one, the target returns the initial value 99 while the source returns 7.
The inequality is checked by the theorem
\begin{center}
\code{oversized_insert_counterexample}.
\end{center}
It witnesses failure of equality for this environment; it does not say
every environment would reveal the defect. If the unwritten contents
happened already to be 7, this particular observation would not distinguish
the translations. Universal quantification over initial contents is what
rules out relying on that accident.

\subsection{Domain rejection and coverage}
Two further theorems prove that $V(0,0)$ and $V(2^{31}-1,2)$ are false.
These are model-domain checks, not executions of the upstream verifier.
Together with the six positive cases and two bad-lowering witnesses,
they constitute ten checked example theorems. They do not measure coverage
of the compiler codebase or prove completeness of the formal grammar.

\subsection{Proof validation and assumptions}
A clean \code{lake build} checked the library and examples. The command
\begin{center}
\code{lake env lean TTIRFormal/Audit.lean}
\end{center}
inspected the assumptions of the principal results. The local lowering
results, universal preservation,
target independence, and both bad-lowering witnesses report
\code{propext} and \code{Quot.sound}. The two integer-bound results
additionally report \code{Classical.choice}. These are standard Lean
foundations, not new compiler axioms introduced by the artifact.

The proof sources contain no \code{sorry}, \code{admit}, custom axiom
declarations, or \code{native_decide} shortcuts. Tactics simplify
equations, discharge finite examples, and solve arithmetic obligations;
the reported evidence is successful proof checking, not a claim that
the specification itself is independently validated. A proof assistant
cannot determine whether a manually encoded denotation is the intended
semantics of a separately implemented compiler.

\section{Source Correspondence and Remaining Obligations}\label{sec:limits}
Table~\ref{tab:source} records the inspected anchors. It is a correspondence
map for audit, not a machine-checked refinement relation. The full revision
is recorded in Section~\ref{sec:source} and in the accompanying findings.

\begin{table}[t]
\caption{Source anchors at the audited revision. Long directory prefixes
are factored into the key below the table.}\label{tab:source}
\centering
\begin{tabularx}{\linewidth}{@{}lX@{}}
\toprule
Anchor & Inspected fact \\
\midrule
A: \code{TritonOps.td}:427 & Splat declaration and tensor result.\\
A: \code{TritonOps.td}:854 & Rank-one i32 range; signed-accessor warning.\\
B: \code{Ops.cpp}:407 & Nonempty interval, length, rank, and element verifier.\\
C: \code{TritonOpConverter.cpp}:1145 & Generated indices and optional start offset.\\
C: \code{TritonOpConverter.cpp}:1205 & Singleton insertion versus fill.\\
C: \code{TritonOpConverter.cpp}:1350 & Broadcast collapse and possible match failure.\\
D: \code{compiler.py}:324,1592 & Stage implementation and registration.\\
\bottomrule
\end{tabularx}
\smallskip
\begin{minipage}{\linewidth}\small
A: \path{include/triton/Dialect/Triton/IR/};
B: \path{lib/Dialect/Triton/IR/};\\
C: \path{third_party/ascend/lib/TritonToLinalg/};
D: \path{third_party/ascend/backend/}.
\end{minipage}
\end{table}

\subsection{What is omitted from the patterns}
The range converter records three attributes:
\begin{itemize}
\item \code{tt.from_make_range};
\item \code{tt.make_range_offset};
\item \code{tt.make_range_size}.
\end{itemize}
The model omits this metadata and does not prove that downstream
transformations interpret it correctly. The splat converter inserts debug
NOP anchors that are likewise omitted. Evaluation models result values,
not every generated operation or diagnostic effect.

We inspected broadcast but did not add it to the formal syntax. Its
converter computes broadcast and unbroadcast dimensions, asks for collapse
reassociation indices, and can fail matching. On success it constructs
collapse-shape and broadcast operations. A proof would need to account
for these shape helpers and success conditions rather than simply adding
an unconditional replication equation. No broadcast theorem is claimed.

\subsection{A layered correctness argument}
The current evidence occupies only part of a potential end-to-end argument:
\begin{enumerate}
\item \emph{Source representation.} Establish that a concrete TTIR fragment
is represented by a well-formed source expression with the same observations.
This requires type, attribute, operand, and relevant layout correspondence.
\item \emph{Abstract lowering.} Prove preservation by $\lo$. This is
Theorem~\ref{thm:main}, completed for the stated grammar.
\item \emph{Implementation correspondence.} Show that the actual converter
implements $\lo$ under explicit assumptions, including index casts,
allocation behavior, and pattern applicability. The present audit is
informal evidence for this step, not a proof.
\item \emph{Surrounding passes and execution.} Establish preservation
through the remaining passes, bufferization, code generation, and execution.
None of these steps is provided by this artifact.
\end{enumerate}
This decomposition explains why a short algebraic preservation proof can
be useful without establishing a large systems claim. The small proof
isolates exactly where the singleton guard and signed offset matter; the
remaining work is to justify that these abstractions describe execution.

\subsection{Threats to validity}
The principal threat is specification error shared by both denotations.
In particular, defining the source range by bit-vector addition makes
its relation to the target transparent, but does not independently
establish the full meaning of the upstream operation. The separate bounds
and concrete signed examples mitigate only part of this issue. A general
signed-decoding theorem and a verified representation relation would
provide stronger evidence.

The pure, total, shape-indexed representation intentionally eliminates
several classes of real compiler behavior: malformed IR, dynamic shapes,
SSA control flow, memory effects, masks, pointer arithmetic, aliasing,
concurrency, and floating point. It also does not distinguish independent
empty allocations or model poison. These omissions limit applicability;
they are not conclusions that the omitted behaviors are safe.

Finally, source inspection is revision-specific and does not automatically
track future changes. The current artifact has no parser or automated
source-to-model validator. Rechecking Lean after a compiler change cannot
by itself establish that the source correspondence still holds.

\section{Related Work}
Triton introduced an intermediate language and compiler for tiled
neural-network computations~\cite{triton}; our artifact is not a
formalization of that entire programming model. It studies two patterns
in a particular downstream checkout. MLIR provides extensible compiler
infrastructure across abstraction levels~\cite{mlir}; our distinction
between a mixed-IR stage and a single dialect follows the practical need
to identify a narrower semantic boundary within such a pipeline. We do
not supply a general framework for composing arbitrary dialect semantics.

CompCert exemplifies machine-checked semantic preservation for a realistic
compiler~\cite{compcert}. The present work uses the familiar structure
of source and target semantics connected by a translation theorem, but
has a far smaller language and lacks a verified connection to the actual
C++ transformation. It should therefore not be read as an alternative
end-to-end verified compiler.

Alive2 addresses bounded translation validation for LLVM~\cite{alive2}.
Our proof is universal over the encoded expression trees rather than
validation of individual LLVM transformations, but the implementation
boundary is weaker: we do not consume before-and-after compiler IR. The
two bad-lowering witnesses are not a translation-validation engine. A
future validator for the selected TTIR/target fragment could complement
the model, provided parsing, unsupported cases, and the validation
semantics were made explicit.

Lean 4 supplies the dependent types, executable functions, and theorem
proving infrastructure used here~\cite{lean4}. We exploit shape indices
to constrain constructors and use function extensionality to reason
about tensors. Our contribution is the source-specific case study,
not a new foundational logic or dependent tensor representation. This
discussion situates the artifact among compiler proofs and validation
approaches; it is not a systematic novelty survey.

\section{Conclusion}
Two elementary-looking tensor conversions expose useful verification
obligations: a singleton write must cover every valid lane, and range
construction must retain the signed offset under finite-width arithmetic.
The Lean core makes both obligations explicit and proves value preservation
under arbitrary input and initial-content families. Its checked
counterexamples distinguish plausible incorrect translations, while
the axiom audit records the logical foundations of the principal results.

The next steps are not larger claims about the same proof. They are
additional evidence: a signed-decoding bridge, a concrete representation
relation, a correspondence argument for the actual conversion patterns,
and preservation through surrounding passes. The present artifact provides
a small verified component and a precise account of what remains outside it.

\section*{Artifact Availability}
The accompanying \code{ttir-formalization} directory contains all Lean
sources, a pinned toolchain, build configuration, and source-audit findings.
From its root, \code{lake build} checks the library and examples;
\code{lake env lean TTIRFormal/Audit.lean} prints the principal axiom
dependencies. The paper resides in its \code{paper} subdirectory.
No public archive DOI or independent artifact-evaluation result is claimed.
The artifact does not require Ascend hardware, and no hardware performance
or compiler-execution experiment is reported.

\section*{Generative AI Disclosure}
Generative AI assisted the development of the accompanying formalization
and the drafting and editing of this manuscript. The reported mechanized
results were checked with Lean. This draft requires author review of the
scientific claims, source correspondence, and references before submission.
AI assistance is not a substitute for that review or for authorship
responsibility.

\bibliography{references}
\end{document}